\section{Controllerdesign}
\label{sec:pilead}

\subsection{PI-Lead design / phase boost, \(\alpha\), \(N_i\)}

\paragraph{Brug når}
Opgaven angiver PI-Lead-form, ønsket crossover \(\omega_c\), phase margin
\(\gamma_M\) og enten \(N_i\) eller en tilsvarende placering af PI-zero.

\paragraph{Input fra opgaven}
Plant \(G(s)\), \(\omega_c\), \(\gamma_M\), \(N_i=\omega_c\tau_i\), og om leadboost
skal placeres maksimalt ved crossover.

\paragraph{Metode}
\begin{enumerate}
\item Skriv controlleren i formen \eqref{eq:pi-lead-controller}. Hvis opgaven giver
      \(N_i=\omega_c\tau_i\), isoler \(\tau_i=N_i/\omega_c\). PI-leddets fase ved
      crossover er \(\phi_{PI}=\angle(1+jN_i)-90^\circ=\arctan(N_i)-90^\circ\).
\item Evaluer plantens fase ved den ønskede crossover: indsæt \(s=j\omega_c\) i
      \(G(s)\), beregn real- og imaginærdel, og brug \(\operatorname{atan2}(\Im,\Re)\).
      Hvis du aflæser fra Bode, brug faseværdien ved netop \(\omega_c\).
\item Skriv faseregnskabet ved crossover:
      \(\phi_G+\phi_{PI}+\phi_{\mathrm{lead}}=-180^\circ+\gamma_M\). Isoler derfor
      \(\phi_{\mathrm{lead}}\). Hvis resultatet er negativt eller urimeligt stort
      (typisk over ca. \(60^\circ\)), passer specifikationen eller strukturen dårligt.
\item Sæt \(\phi_{\max}=\phi_{\mathrm{lead}}\), når leadmaksimum skal ligge ved
      \(\omega_c\). Beregn \(\alpha\) og \(\tau_d\) med
      \eqref{eq:pi-lead-parameters}; kontrollér straks at \(0<\alpha<1\) og
      \(\tau_d>0\).
\item Bestem \(K_P\) ved at fjerne \(K_P\) fra controlleren, evaluere
      \(C_0(j\omega_c)G(j\omega_c)\), og bruge
      \(K_P=1/|C_0(j\omega_c)G(j\omega_c)|\) fra \eqref{eq:pi-lead-magnitude}.
\item Kontrollér til sidst at det færdige loop har magnitude \(1\) og fase
      \(-180^\circ+\gamma_M\) ved \(\omega_c\); ellers er der en enheds-, grad/radian-
      eller fortegnsfejl.
\end{enumerate}

\paragraph{Formler}
\begin{equation}
\label{eq:pi-lead-controller}
C(s)=K_P\frac{\tau_i s+1}{\tau_i s}
          \frac{\tau_d s+1}{\alpha\tau_d s+1},\qquad 0<\alpha<1,
\end{equation}
\begin{equation}
\label{eq:pi-lead-parameters}
\phi_{\max}=\sin^{-1}\!\left(\frac{1-\alpha}{1+\alpha}\right),\quad
\alpha=\frac{1-\sin\phi_{\max}}{1+\sin\phi_{\max}},\quad
\tau_d=\frac1{\omega_c\sqrt{\alpha}},
\end{equation}
\begin{equation}
\label{eq:pi-lead-magnitude}
\left|C(\mathrm{j}\omega_c)G(\mathrm{j}\omega_c)\right|=1 .
\end{equation}

\paragraph{Symboler}
\(K_P\) er proportionalgain; \(\tau_i,\tau_d\) [s] er tidskonstanter;
\(\alpha\) og \(N_i\) er dimensionsløse.

\paragraph{Antagelser}
PI anvendes fordi integratorens stationære fordel ønskes; leadboostets maximum
placeres ved ønsket crossover; plantmodellen er gyldig omkring operationen.

\paragraph{Hurtigtjek}
\(0<\alpha<1\), positive tidskonstanter og gain; det beregnede loop skal have
magnitude 1 og fase \(-180^\circ+\gamma_M\) ved \(\omega_c\).
For et rent P-Lead-led uden PI bruges samme leaddel:
\[
C_D(s)=\frac{\tau_d s+1}{\alpha\tau_d s+1},\qquad
\omega_{\max}=\frac{1}{\tau_d\sqrt{\alpha}},\qquad
|C_D(j\omega_{\max})|=\frac1{\sqrt{\alpha}}.
\]
Hvis \(\omega_c\) og \(\alpha\) er givet, er
\(\tau_d=1/(\omega_c\sqrt{\alpha})\), lead-zero-frekvensen er \(1/\tau_d\), og
lead-polfrekvensen er \(1/(\alpha\tau_d)\). I dB er leadbidraget ved centerfrekvens
\(20\log_{10}(1/\sqrt{\alpha})\).

\paragraph{Typiske fælder}
At glemme den negative PI-fase eller at bruge grader som radianer. F21 Q18's
solutiontekst og direkte TF-evaluering afviger lidt i plantfasen, men begge vælger
\(\alpha\approx0.08\), \(K_P\approx200\). For P-Lead: at bytte zero og pole eller
svare den lineære faktor \(1/\sqrt{\alpha}\) som om den var dB.

\paragraph{Python}
\code{design\_pi\_lead\_at\_crossover(num,den,omega\_c,phase\_margin\_deg,n\_i)}
udfører netop ovenstående beregning og returnerer targetchecks. Vælg ikke controllerstruktur
eller crossover blindt ud fra outputtet.

\subsection{P-gain fra margin / proportional controller design}

\paragraph{Brug når}
Der skal findes et \(K_P\) ud fra en ønsket phase margin eller stabilitetsgrænse.

\paragraph{Input fra opgaven}
Plantens magnitude/fase eller analytiske \(G(s)\), samt ønsket margin.

\paragraph{Metode}
\begin{enumerate}
\item Positivt P-gain ændrer kun magnitude, ikke fase. Find derfor den frekvens
      \(\omega_c\), hvor plantfasen alene opfylder første del af
      \eqref{eq:p-gain-margin-design}: \(\angle G(j\omega_c)=-180^\circ+\gamma_M\).
\item Hvis du har en formel for \(G(s)\), indsæt \(s=j\omega\), brug
      \(\operatorname{atan2}\) for fasen, og løs for \(\omega\). Hvis du har Bode-plot,
      aflæs den frekvens hvor fasekurven rammer målfasen.
\item Aflæs eller beregn magnituden \(M=|G(j\omega_c)|\). Er den angivet i dB, konvertér
      først med \eqref{eq:bode-db-conversion}.
\item Sæt \(K_P=1/M\) med \eqref{eq:p-gain-margin-design}, så
      \(|K_PG(j\omega_c)|=1\). Ved stabilitetsgrænse/gain margin skal marginalpunktet
      ikke vælges som stabilt; vælg et gain inde i det stabile interval.
\item Kontrollér closed-loop-polerne for det valgte \(K_P\). Hvis opgaven også spørger
      om stationær fejl, brug \eqref{eq:steady-state-error} efter poltesten.
\end{enumerate}

\paragraph{Formler}
\begin{equation}
\label{eq:p-gain-margin-design}
\angle G(\mathrm{j}\omega_c)=-180^\circ+\gamma_M,\qquad
K_P=\frac1{|G(\mathrm{j}\omega_c)|}.
\end{equation}

\paragraph{Symboler}
\(K_P\) er dimensionsløs gain i opgavernes normalisering; \(\gamma_M\) måles i grader.

\paragraph{Antagelser}
Positiv P-gain påvirker kun magnitude, og plantens stabilitetsforudsætninger skal
stemme med det anvendte Bodeargument.

\paragraph{Hurtigtjek}
F21 Q6: \(38.9\) dB kræver \(K_P\approx88\), ikke 38.9.
Efter et design skal hele Bode/Nyquist-kurven kontrolleres. Et loop med resonans kan
have et ekstra magnitude peak over 0 dB ved fase under \(-180^\circ\), selv om det
ønskede crossoverpunkt er designet korrekt.

\paragraph{Typiske fælder}
At acceptere marginal gain som stabilt eller at overse type-0 stationær fejl.

\paragraph{Python}
\code{evaluate\_transfer\_function} og \code{closed\_loop\_poles} kan kontrollere
en analytisk beregning, men kan ikke aflæse Bodefiguren.

\subsection{Ziegler--Nichols PID fra ultimativ gain / \(K_u\), \(P_u\)}
\label{sec:ziegler}

\paragraph{Brug når}
Opgaven giver en PID-loopstruktur, en kritisk/ultimativ proportionalgain \(K_u\), eller
beder om Ziegler--Nichols tuning fra sustained oscillation.

\paragraph{Input fra opgaven}
Kritisk gain \(K_u\), oscillationsfrekvens \(\omega_u\) eller periode \(P_u\), og at
controlleren skal skrives som ideal/parallel PID-form.

\paragraph{Metode}
\begin{enumerate}
\item Hvis \(K_u\) ikke er opgivet, find marginalstabilitet for P-only closed loop:
      skriv den karakteristiske nævner, sæt Routh-grænsen eller \(s=j\omega\), og løs
      for \(K_u\) og \(\omega_u\).
\item Beregn perioden med \eqref{eq:ultimate-period}. Pas på brøken: perioden er
      \(2\pi/\omega_u\), ikke \(\omega_u/(2\pi)\).
\item Brug PID-rækken i \eqref{eq:ziegler-nichols-pid}: \(K_P=0.6K_u\),
      \(T_i=P_u/2\), \(T_d=P_u/8\).
\item Indsæt tallene i \eqref{eq:ideal-pid}. Hvis svarmulighederne har formen
      \(K_P(1+1/(T_i s)+T_d s)\), skal både gainet udenfor parentesen og tiderne inde i
      parentesen matche.
\end{enumerate}

\paragraph{Formler}
\begin{equation}
\label{eq:ultimate-period}
P_u=\frac{2\pi}{\omega_u}.
\end{equation}
\begin{equation}
\label{eq:ziegler-nichols-pid}
K_P=0.6K_u,\qquad T_i=\frac{P_u}{2},\qquad T_d=\frac{P_u}{8}.
\end{equation}
\begin{equation}
\label{eq:ideal-pid}
C_{\mathrm{PID}}(s)=K_P\left(1+\frac{1}{T_i s}+T_d s\right).
\end{equation}

\paragraph{Symboler}
\(K_u\) er ultimativ gain; \(P_u\) er ultimativ periode; \(T_i,T_d\) er integral- og
derivative-tid.

\paragraph{Antagelser}
Reglen er en heuristisk tuningregel for P-only marginal oscillation; den garanterer ikke
god robusthed eller acceptable actuatorlimits.

\paragraph{Hurtigtjek}
F25 Q8: \(K_u=30\), \(\omega_u=\sqrt5\), \(P_u=2.809\), så \(K_P=18\),
\(T_i=1.405\), \(T_d=0.351\).

\paragraph{Typiske fælder}
At bruge \(K_u\) direkte som \(K_P\), at bytte \(P_u/2\) og \(P_u/8\), eller at vende
\(2\pi/\omega_u\).

\subsection{Cascaded control og REGBOT balance / inner and outer loops}
\label{sec:cascade-regbot}

\paragraph{Brug når}
Slides omtaler REGBOT balance, inverted pendulum, tilt-controller, velocity-controller
eller position-controller i kaskade.

\paragraph{Metode}
\begin{enumerate}
\item Design indefra og ud. Det inderste loop skal være stabilt og hurtigere end det
      loop, der bruger det som aktuatormodel.
\item For REGBOT balance er den fysiske tolkning: motor/wheel dynamics giver kraft eller
      acceleration, tilt-loopet stabiliserer den ustabile pendulvinkel, velocity-loopet
      laver en tilt-reference, og position-loopet laver en velocity-reference.
\item Når et indre loop er lukket, erstattes det i ydre design af dets closed-loop
      transferfunktion eller en lavordens approximation, ikke af den oprindelige åbne
      plant.
\item Vælg ydre crossover tydeligt lavere end indre crossover, så tidskonstanterne er
      separerede. Hvis loops ligger for tæt, ødelægger de hinandens phase margin.
\item Kontroller til sidst hele kaskaden med actuatorlimits, fordi et stabilt lineært
      inderloop stadig kan kræve for stor motorhastighed, acceleration eller tilt-vinkel.
\end{enumerate}

\paragraph{Hurtigtjek}
Et open-loop-ustabilt balanceled skal først stabiliseres; position eller velocity kan
ikke designes robust oven på et ustabilt tilt-loop. I en eksamensforklaring er det ofte
nok at angive looprækkefølgen, separationen i bandwidth og hvorfor referencesignalet
går fra ydre til indre loop.

\paragraph{Typiske fælder}
At tune position-loopet direkte på den ustabile pendulplant; at gøre alle loops lige
hurtige; at validere hvert loop isoleret uden at kontrollere den samlede kaskade.
